\subsection{Cancel -- Cancellations}

The \texttt{cancel} package is used to visually cancel terms in
mathematical expressions.


\subsubsection{Basic Cancellation}

\begin{lstlisting}[style=latex]
\[
\cancel{x}
\]
\end{lstlisting}

\[
\cancel{x}
\]


\subsubsection{Cancelling Common Terms}

It can be used to show the simplification of common terms:

\begin{lstlisting}[style=latex]
\[
\frac{\cancel{x}y}{\cancel{x}z}
=
\frac{y}{z}
\]
\end{lstlisting}

\[
\frac{\cancel{x}y}{\cancel{x}z}
=
\frac{y}{z}
\]


\subsubsection{Cancelling to a Value}

The \verb|\cancelto| command indicates the value to which a term
is reduced:

\begin{lstlisting}[style=latex]
\[
a + \cancelto{0}{b-b} = a
\]
\end{lstlisting}

\[
a + \cancelto{0}{b-b} = a
\]


\subsubsection{Other Cancellation Styles}

Different cancellation directions are also available:

\begin{lstlisting}[style=latex]
\[
\cancel{x}
\qquad
\bcancel{x}
\qquad
\xcancel{x}
\]
\end{lstlisting}

\[
\cancel{x}
\qquad
\bcancel{x}
\qquad
\xcancel{x}
\]